% ===== §10 =====
\section{Legitimacy without Correctness}
\label{sec:legitimacy}

\noindent
The distinction the functional case cannot draw is a normative one, and it must be drawn with care, since the obvious way of drawing it is closed to this account. The obvious way would be to say that genuine understanding is correct where the counterfeit is mistaken, that the understanding partner grasps the other as the other truly is while the managing partner works from a false or partial picture, so that the distinction between understanding and extraction is the distinction between a true and a false representation of the other. This way is closed, because the epistemology on which the paper stands does not grant that an understanding of another is ever correct in the sense this would require.

The reason returns to the account of truth recalled in Part One. Understanding another, on this account, is a generative observation, produced in the crossing of two worlds, historically situated, and always incomplete. It is never the possession of the other as the other finally is, for there is no rendering that holds the other without remainder, and the observation of the other is enriched without end as the crossing continues and the two worlds change. Truth, here as everywhere on this account, is generated in a historical and contradictory movement and is never held as a finished correctness. To defend genuine understanding as the correct representation of the other, against the counterfeit as the false one, would be to claim for understanding exactly the finished correctness the epistemology denies to all knowing. The account cannot draw its normative distinction as a distinction in correctness without abandoning its own conception of truth.

The distinction must therefore be drawn on other ground, and the ground is ethical and not epistemic. What separates understanding from its counterfeit is not that the one is correct and the other mistaken. It is that the one is legitimate and the other is not, and the legitimacy in question is a matter of how the other is treated in the crossing, not of whether the crossing arrives at a correct result. The understanding partner enters the other's world as another world, to be observed with and not possessed, and holds the other as a co-generating subject whose rendering of things is a source and not a resource. The managing partner enters the other's world as a terrain to be mapped for the sake of a hold, and holds the other as an object of management, a system to be kept in a desired state. The difference is in the standing granted to the other, and it is a difference of legitimacy.

\begin{claim}[Legitimacy, not correctness]
The account does not distinguish genuine understanding from its counterfeit by the correctness of the one and the error of the other, for understanding another is a generative observation, historically produced and always incomplete, and is never held as a finished correctness. The distinction is ethical. Genuine understanding is legitimate in that it enters the other's world to observe with the other and holds the other as a co-generating subject; its counterfeit is illegitimate in that it enters to secure a hold and holds the other as an object of management. What is asked of an understanding is that it be legitimate and not that it be correct.
\end{claim}

The ground of this legitimacy has been carried, in the wider framework, by an ethic of non-possession, and non-possession will remain central to what follows. To understand another legitimately is to observe the other without possessing the other, to let the other's world be a world one enters and does not annex, and to hold what one observes of the other as a matter generated between the two and belonging to neither, and not as a property of one's own by which the other is held. The counterfeit fails this at its root, for its whole motion is possessive, a taking of the other into a system of management under the appearance of a giving of attention. This is why the framework names it a relational extraction, and why the wrong it does is legible in the terms of generative justice, as the interruption of the mutual generation that a legitimate bond sustains, the rendering of one party into fuel for the other's ends under the form of care.

Non-possession states one aspect of the legitimacy in question, and the account should not rest the whole of its normative ground on a single ethical rendering, for a reason internal to its own method. Legitimacy is a concept, a matter to be understood, and the epistemology this paper applies holds that no single symbolic system exhausts a concept and that a concept becomes observable through the relations among several renderings of it. An ethical tradition is such a symbolic system. To render the legitimacy of understanding through non-possession alone would repeat, at the level of ethics, the error the epistemology sets itself against, the mistaking of one rendering for the whole. The next section therefore observes ethical legitimacy across several traditions, of which non-possession names one aspect, so that the normative ground is drawn in the manner the paper's own method requires.

Two further points hold this normative ground to the conception of truth the paper shares, so that the normative case does not smuggle back, at the level of ethics, the finished correctness it refused at the level of knowing. The first is that legitimacy, as it is used here, makes no claim to have grasped the other correctly. A legitimate understanding may be deeply incomplete, may err, may be revised across the history of the bond, and remains legitimate throughout, since its legitimacy lies in the standing it grants the other and not in the accuracy of its result. One may hold one's partner as a co-generating subject, enter their world with the intent to observe and not to possess, and still understand them partially and by degrees, as the account holds all understanding to be. Legitimacy is a matter of the ethical form of the crossing, compatible at every point with the incompleteness the epistemology requires.

The second is that the criterion of legitimacy is itself historically generated, and is set here as no final guarantee. The ethic of non-possession, the standing of the other as a co-generating subject, the very conception of a relational extraction, are not timeless moral facts read off from the nature of things. They are themselves generated in the historical and material development of relational life, arising as the conditions of a form of life bring them forth and open to the same dialectical movement as everything else the account treats. What counts as possession and what as legitimate observation of another, where the line falls between entering a world and annexing it, is a matter that has a history and will have a further history, refigured as the conditions of intimate life develop. The account does not install non-possession as a new absolute in the place of the correctness it declined. It offers non-possession as the criterion legitimacy takes in the present development of relational life, held as the best the present affords and answerable to its own further generation, in exactly the manner the framework holds of all its commitments, including this one.

With this the distinction is drawn in principle. Function tells what understanding does and cannot tell understanding from its counterfeit. Legitimacy tells them apart, without claiming the correctness the account denies to all knowing, and holding even its own criterion within the history it theorizes. What remains is to draw the ethical ground of legitimacy more fully than a single rendering allows, by observing it across several ethical traditions, and this is the work of the section that follows.
